\section{A preservation theorem for compositions}\label{sec:preservation}
This section carries out steps~(i) and~(ii) of the outline in
Section~\ref{ov:outline}: it builds the structure $\M$ and proves the
preservation theorem (Theorem~\ref{thm:preservation}), for an arbitrary
kernel in place of conjunction. The construction does not mention $\Psi_k$
or $\Phi_k$, whose interpretation is step~(iii) (Section~\ref{sec:logical}).
The proof uses just one property of composition, isolated in
Section~\ref{sec:responses}, and this is what makes the theorem applicable
to matrix products in Section~\ref{sec:matrix}.

\paragraph{The shape of the argument.}
The argument has three independent parameters, and each of the sections
that follow varies one of them while holding the other two fixed: the
interpretation of the source query on $\M$ (Section~\ref{sec:logical}, for
$\Psi_k$ and for $\Phi_k$), the family of perturbations of $\M$ (deletion
of cells here, recoloring of the hub's edges on a fixed domain in
Section~\ref{sec:robustness}), and the operation whose occurrences are
counted (conjunction in Section~\ref{sec:responses}, a semiring product and
an arbitrary kernel in Section~\ref{sec:matrix}). A single number, the
support number of Definition~\ref{def:width}, couples the three, and it
depends only on the operation and the join. Everything else is proved here
for an arbitrary kernel, so Section~\ref{sec:matrix} only has to recompute
$\wid\odot$; this is why the construction that separates \FOthree\ from
\CoR\ also answers the question of \citet{brijder2019} about MATLANG
(Section~\ref{semi:typed-section}).

Two features of the conclusion are worth recording, since they are not
usually available together. First, the bound is proved against a fixed
finite family of $|\Col|+1$ structures, chosen in advance of the candidate
and written down explicitly (Proposition~\ref{prop:explicit}); no limit is
taken and no infinite family is involved. Second, a shared subterm is
charged only once, so the statement concerns circuits and not only terms.
We claim no more than this: the invariant separates one structure from a
bounded family of its perturbations, and the bound accordingly counts
occurrences of a single operation in a language whose remaining operations
are free. It is not a lower bound for Boolean circuits with access to
individual entries of the input.

\subsection{The structure}\label{sec:structure}
Fix a finite set $\Col$ of \emph{colors} with $N=|\Col|\ge2$. The vertices
of $\M$ are of three kinds (Figure~\ref{fig:structure} on
page~\pageref{fig:structure}):
\begin{itemize}[nosep]
\item one \emph{hub} $h$;
\item for each color $\lambda\in\Col$ a \emph{cell} $\cell\lambda$ of
 $L=\lceil12N^2\ln(2N)\rceil$ \emph{cell vertices}; the cells are disjoint,
 and $\cell\Col=\bigcup_{\lambda\in\Col}\cell\lambda$;
\item a finite set $X$ of \emph{anchors}, one per color in
 Section~\ref{sec:logical} and possibly more in
 Section~\ref{sec:robustness}.
\end{itemize}
The \emph{core} is $U=\{h\}\cup\cell\Col$ and the domain is $D=X\cup U$. A
pair of \emph{distinct} vertices is called an \emph{edge}. Every edge
$\{u,v\}$ between two core vertices carries a color
$\colr(u,v)=\colr(v,u)\in\Col$; edges incident to an anchor carry none. The
hub's edges announce the cell,
\begin{equation}\label{eq:hub-colours}
 \colr(h,u)=\lambda\qquad\text{for }u\in\cell\lambda,
\end{equation}
and the edges between cell vertices are colored so that the following holds.

\begin{lemma}[Extension property]\label{lem:palette}
The edges between cell vertices can be colored so that
\begin{enumerate}[nosep,label=(\alph*)]
\item for all distinct $u,v\in\cell\Col$, all colors $\mu,\nu\in\Col$, and
 every cell $\cell\lambda$, some $z\in\cell\lambda\setminus\{u,v\}$
 satisfies $\colr(u,z)=\mu$ and $\colr(z,v)=\nu$;
\item for every $u\in\cell\Col$, every color $\mu\in\Col$, and every cell
 $\cell\lambda$, some $z\in\cell\lambda\setminus\{u\}$ satisfies
 $\colr(u,z)=\mu$.
\end{enumerate}
\end{lemma}
\begin{proof}
Part (b) follows from (a) by choosing any $v\ne u$ and any $\nu$. For (a),
color the edges between cell vertices independently and uniformly at random.
A \emph{requirement} is a choice of $\lambda$, of distinct
$u,v\in\cell\Col$, and of $\mu,\nu$; since $|\cell\Col|=NL$ there are fewer
than $N\cdot(NL)^2\cdot N^2=N^5L^2$ of them. For a fixed requirement, each of
the at least $L-2$ vertices $z\in\cell\lambda\setminus\{u,v\}$ satisfies it
with probability $N^{-2}$, and these events are independent because they
concern disjoint pairs of edges, so the requirement fails with probability at
most $(1-N^{-2})^{L-2}\le\exp(-(L-2)/N^2)$. As $\ln(2N)\le N$ and $N\ge2$ we
have $L\le13N^3$ and $(L-2)/N^2\ge12\ln(2N)-\tfrac12$, so the probability
that some requirement fails is at most
\[
 N^5L^2\exp(-(L-2)/N^2)
 \ \le\ 169N^{11}\exp(1/2)\,(2N)^{-12}
 \ =\ \frac{169\exp(1/2)}{4096\,N}\ <\ 1 .
\]
Hence a coloring with property (a) exists.
\end{proof}
The constant $12$ is sufficient rather than optimal: it is what produces
the factor $1/N$ in the bound just displayed. Nothing below depends on its
value, since $L$ enters the results of this paper only through the size
threshold \eqref{eq:size-threshold}.

We fix such a coloring once and for all. As explained in
Section~\ref{ov:structure}, only the hub distinguishes cells, by
\eqref{eq:hub-colours} and Lemma~\ref{lem:palette}, and the anchors, which
are never deleted, will be distinguished by the interpretation of the
relation symbols. A pair $(x,h)$ with $x$ an anchor is called a
\emph{read-off pair}; these are the only pairs at which
Section~\ref{sec:logical} evaluates the two hard queries.

The coloring of Lemma~\ref{lem:palette} is the only object in this paper
obtained from a probabilistic existence argument, and this argument can be
avoided: Appendix~\ref{app:explicit} gives a deterministic polynomial-time
construction and a closed form over a prime field, in which the cells are
the cosets of the $N$-th powers (Proposition~\ref{prop:explicit}). Either
construction makes the $|\Col|+1$ structures of
Corollary~\ref{cor:detectable} explicit, and with the closed form the size
threshold \eqref{eq:size-threshold} remains polynomial in $N$. The cells
cannot be made much smaller: property~(a) of Lemma~\ref{lem:palette}
already forces $L\ge N^2+1$ (Remark~\ref{rem:covering}).

\subsection{Types, homogeneity, and profiles}\label{sec:types}
Every ordered pair of vertices receives a \emph{type}, by cases:
\begin{equation}\label{eq:types}
 \tp(u,v)=
 \begin{cases}
  \tpa{x}{x'}&\text{if }u=x\in X\text{ and }v=x'\in X,\\
  \tpl{x}&\text{if }u=x\in X\text{ and }v\in U,\\
  \tpr{x}&\text{if }u\in U\text{ and }v=x\in X,\\
  \tpeq&\text{if }u=v\in U,\\
  \ctype{\colr(u,v)}&\text{if }u\ne v\text{ and }u,v\in U.
 \end{cases}
\end{equation}
As sets of pairs, $\tpa{x}{x'}=\{(x,x')\}$, $\tpl{x}=\{x\}\times U$,
$\tpr{x}=U\times\{x\}$, $\tpeq=\Delta_U$, and
$\ctype\lambda=\{(u,v)\in U^2:u\ne v,\ \colr(u,v)=\lambda\}$; they partition
$D^2$. There are $|X|^2+2|X|+1+|\Col|$ types, however large the cells are.

Note that a cell is not a type, and neither is $\{h\}$: the pairs $(h,u)$
with $u\in\cell\lambda$ lie in $\ctype\lambda$ together with every other
edge of color $\lambda$ (Section~\ref{ov:types}).

An array $P:D^2\to S$ with values in a set $S$ --- in particular a relation,
where $S$ is the set $\bool=\{0,1\}$ of truth values --- is
\emph{type-constant} if it is constant on every type;
we then write $P(\tau)$ for its value there. The point of this subsection is
that type-constancy propagates to everything computed from type-constant
inputs (Section~\ref{ov:types}), because a witness is seen only through the
types of its two edges: for type-constant $P,Q$, the term
$P(u,z)\land Q(z,v)$ that a witness $z$ contributes to $(P;Q)(u,v)$ depends
only on the \emph{profile} $\bigl(\tp(u,z),\tp(z,v)\bigr)$ of $z$ at
$(u,v)$, so the value of the composition is determined by the set of
profiles that some witness realizes. The structure is built so that this
set depends only on the type of the pair.

\begin{lemma}[Homogeneity]\label{lem:profile}
For every type $\tau$, the set of witness profiles
\[
 W_\tau=\{(\tp(u,z),\tp(z,v)):z\in D\}
\]
is the same for all pairs $(u,v)\in\tau$.
\end{lemma}
\begin{proof}
Let $(u,v)\in\tau$. A witness $z$ is an anchor, one of the two endpoints, or
another core vertex; we list the profiles of each kind and check that they
do not depend on the choice of $(u,v)$ in $\tau$.

\emph{Case $\tau=\ctype\lambda$.} Then $u\ne v$ are core vertices with
$\colr(u,v)=\lambda$. An anchor $z=x$ gives $(\tpr{x},\tpl{x})$; $z=u$ gives
$(\tpeq,\ctype\lambda)$ and $z=v$ gives $(\ctype\lambda,\tpeq)$. A core
witness $z\notin\{u,v\}$ gives $(\ctype{\mu},\ctype{\nu})$ with
$\mu=\colr(u,z)$ and $\nu=\colr(z,v)$, and every pair $(\mu,\nu)$ occurs:
for $u,v\in\cell\Col$ by Lemma~\ref{lem:palette}(a); for $u=h$ choose
$z\in\cell\mu\setminus\{v\}$ with $\colr(z,v)=\nu$ by
Lemma~\ref{lem:palette}(b), so that $\colr(h,z)=\mu$ by
\eqref{eq:hub-colours}; the case $v=h$ is symmetric.

\emph{Case $\tau=\tpeq$.} Then $u=v\in U$. Anchors contribute as before, $z=u$
gives $(\tpeq,\tpeq)$, and a core witness $z\ne u$ gives
$(\ctype\mu,\ctype\mu)$ with $\mu=\colr(u,z)$; every $\mu$ occurs, by
Lemma~\ref{lem:palette}(b) if $u\in\cell\Col$ and by
\eqref{eq:hub-colours} if $u=h$.

\emph{Case $\tau=\tpl{x}$.} Then $u=x$ is an anchor and $v\in U$. An anchor $z=x'$
(possibly $x'=x$) gives $(\tpa{x}{x'},\tpl{x'})$; $z=v$ gives
$(\tpl{x},\tpeq)$; a core witness $z\ne v$ gives $(\tpl{x},\ctype\mu)$ with
$\mu=\colr(z,v)$, and every $\mu$ occurs as in the previous case. The type
$\tpr{x}$ is symmetric, and a type $\tpa{x}{x'}$ consists of a single pair.
\end{proof}

Homogeneity is exactly the property that composition needs.
\begin{corollary}[Type-constant relations are closed]\label{cor:closure}
Let $P,Q$ be type-constant relations on $D$. Then $P;Q$, $P^\smile$,
$\overline P$, $P\cup Q$, $P\cap Q$, and the constants $\zero,\one,\id$ are
type-constant, and for every type $\tau$
\begin{equation}\label{eq:profile-value}
 (P;Q)(\tau)=\bigvee_{(\tau_1,\tau_2)\in W_\tau}P(\tau_1)\land Q(\tau_2).
\end{equation}
\end{corollary}
\begin{proof}
For $(u,v)\in\tau$ we have
$(P;Q)(u,v)=\bigvee_{z\in D}P(\tp(u,z))\land Q(\tp(z,v))$, and by
Lemma~\ref{lem:profile} the set of profiles over which the join ranges
depends only on $\tau$; this is \eqref{eq:profile-value}. Union,
intersection, and complement act entrywise. Converse maps each type onto a
type ($\ctype\lambda$ and $\tpeq$ onto themselves, $\tpl{x}$ onto $\tpr{x}$,
$\tpa{x}{x'}$ onto $\tpa{x'}{x}$). The constants $\zero$ and $\one$ are constant,
and $\id$ is true exactly on the types $\tpeq$ and $\tpa{x}{x}$.
\end{proof}
Thus every relation computed on $\M$ from type-constant inputs is a finite
table with one entry per type.

\subsection{Deleting cells: three responses}\label{sec:responses}
For a nonempty set $J\subseteq\Col$ let
\[
 D_J=X\cup\{h\}\cup\cell J,\qquad \cell J=\bigcup_{\lambda\in J}\cell\lambda,
\]
and write $\M\res J$ and $P\res J$ for $\M\res D_J$ and $P\res D_J$: the
cells of all colors outside $J$ are deleted and everything else is
retained. We ask how the value of a composition on $\M\res J$ depends on
$J$.

\paragraph{Compositions as aggregations.}
It costs nothing to answer this question for a more general operation,
which is what Section~\ref{sec:matrix} needs. Let $(S,\vee,0)$ be a join-semilattice with
least element $0$, so that finite joins are associative, commutative and
idempotent, and let $\odot$ be any binary operation on $S$, called the
\emph{kernel}. For arrays $P,Q:D'^2\to S$ on a domain $D'$, the
\emph{aggregation} of $P$ and $Q$ with kernel $\odot$ is
\begin{equation}\label{eq:aggregation}
 (P\agg\odot Q)(u,v)=\bigvee_{z\in D'}P(u,z)\odot Q(z,v).
\end{equation}
For $S=\bool=\{0,1\}$, with disjunction as the join and $\odot=\land$, this
is composition,
$P\agg\land Q=P;Q$, and a reader interested only in
Theorems~\ref{main:omega} and~\ref{main:phi} may read $p\odot q$ as
$p\land q$ throughout. No algebraic law is assumed for $\odot$. A
\emph{computation} is a circuit whose gates are aggregations, possibly with
different kernels at different gates, together with arbitrary functions
$S^r\to S$ applied entry by entry, transpose, constant arrays, and the
equality array, which takes two fixed distinct values on and off the
diagonal. For $S=\bool$ this includes every \CoR\ term and every relational
circuit. The gates other than aggregations are \emph{free}: they commute
with restriction to a subdomain, and they preserve type-constancy.
Lemma~\ref{lem:profile} and \eqref{eq:profile-value} hold for every kernel,
with $\odot$ in place of $\land$, since only the idempotence of $\vee$ was
used in Corollary~\ref{cor:closure}.

\paragraph{Where the retained colors enter.}
Let $P,Q$ be type-constant arrays on $D$ and write
\[
 p_\lambda=P(\ctype\lambda),\qquad q_\lambda=Q(\ctype\lambda)
 \qquad(\lambda\in\Col)
\]
for their values on the edges of color $\lambda$. As explained in
Section~\ref{ov:twocolors}, whose argument does not depend on the kernel, a
deletion can change the value of an aggregation only at a pair with a hub
endpoint, and the hub sees the whole of $\cell\lambda$ through $p_\lambda$
and $q_\lambda$ alone: a retained witness $z\in\cell\lambda$ contributes a
term $p_\lambda\odot s$, $s\odot q_\lambda$ or $p_\lambda\odot q_\lambda$,
where $s$ is the value carried by its other edge, according to whether the
hub is the left endpoint, the right endpoint, or both. This motivates the
three \emph{responses} of the family $(p_\lambda,q_\lambda)_{\lambda\in\Col}$
to a set $J$ of colors:
\begin{equation}\label{eq:responses}
 \ell_J(s)=\bigvee_{\lambda\in J}p_\lambda\odot s,\qquad
 r_J(s)=\bigvee_{\lambda\in J}s\odot q_\lambda,\qquad
 d_J=\bigvee_{\lambda\in J}p_\lambda\odot q_\lambda
 \qquad(s\in S).
\end{equation}
The \emph{left} and \emph{right responses} $\ell_J$ and $r_J$ are functions
on $S$; the \emph{paired response} $d_J$ is a scalar. For $\odot=\land$ they amount to the three bits of Section~\ref{ov:twocolors}.

These three responses are all that a deletion can disturb.
\begin{lemma}[Three responses preserve an aggregation]\label{lem:local}
Let $P,Q$ be type-constant arrays on $D$ and $\varnothing\ne J\subseteq\Col$.
If $(\ell_J,r_J,d_J)=(\ell_{\Col},r_{\Col},d_{\Col})$, then
\[
 (P\agg\odot Q)\res J=(P\res J)\agg\odot(Q\res J),
\]
where the right-hand side is evaluated on $D_J$.
\end{lemma}
\begin{proof}
Let $u,v\in D_J$. Both sides are joins over witnesses, the left one over
$z\in D$ and the right one over $z\in D_J$. Witnesses in $X\cup\{h\}$ and
witnesses in $\{u,v\}$ lie in $D_J$ and contribute the same term to both
sides. It remains to compare, for $I=J$ and for $I=\Col$, the contribution
of the cell witnesses,
\[
 \sigma_I(u,v)=\bigvee_{z\in\cell I\setminus\{u,v\}}
 P(u,z)\odot Q(z,v).
\]
As $P$ and $Q$ are type-constant, $\sigma_I(u,v)$ is determined by the set
of profiles realized by $z\in\cell I\setminus\{u,v\}$. We read these off
from \eqref{eq:hub-colours} and Lemma~\ref{lem:palette} for each kind of
endpoint. Following the rule $p_\lambda=P(\ctype\lambda)$, write $p_x=P(\tpl{x})$ and
$q_{x'}=Q(\tpr{x'})$; the two subscript positions differ because $P$ is
always the left factor of the aggregation and $Q$ the right one. The result
is
\begin{equation}\label{eq:response-table}
\renewcommand{\arraystretch}{1.35}
\begin{array}{c|ccc}
 \sigma_I(u,v)& v=x'\in X&v=h&v\in\cell J\\ \hline
 u=x\in X&p_x\odot q_{x'}&r_I(p_x)&r_{\Col}(p_x)\\[1pt]
 u=h&\ell_I(q_{x'})&d_I&\displaystyle\bigvee_{\mu\in\Col}\ell_I(q_\mu)\\[3pt]
 u\in\cell J&\ell_{\Col}(q_{x'})&\displaystyle\bigvee_{\mu\in\Col}r_I(p_\mu)&
 \begin{cases}d_{\Col}&\text{if }u=v,\\[1pt]
 \displaystyle\bigvee_{\mu,\nu\in\Col}p_\mu\odot q_\nu&\text{if }u\ne v.\end{cases}
\end{array}
\end{equation}
We verify the table row by row. For two anchor endpoints, every cell witness
contributes the same term $p_x\odot q_{x'}$, so the entry does not depend on
$I$. For an anchor $x$ and the hub, a witness $z\in\cell\lambda$ has
$\tp(z,h)=\ctype\lambda$, so the retained colors $\lambda\in I$ enter
through $r_I$. For an anchor
and a cell vertex $v$, the edges $\{z,v\}$ with $z$ in any single retained
cell already realize every color, by Lemma~\ref{lem:palette}(b), so the
entry does not depend on $I$. In the second row a witness $z\in\cell\lambda$
is seen from the hub through $p_\lambda$: at $(h,x')$ this gives
$\ell_I(q_{x'})$; at $(h,h)$ the two edges of $z$ have the same color
$\lambda$, which gives $d_I$; at $(h,v)$ with $v$ a cell vertex the profiles
are $(\ctype\lambda,\ctype\mu)$ with $\lambda\in I$ and $\mu$ arbitrary, by
Lemma~\ref{lem:palette}(b). The first two entries of the third row are
obtained from the third column by exchanging the roles of the two endpoints.
Finally, for two cell-vertex endpoints, Lemma~\ref{lem:palette}(a) yields
every ordered pair of colors if $u\ne v$, while for $u=v$ both edges of $z$
have the same color and every color occurs.

Only the entries containing $\ell_I$, $r_I$ or $d_I$ depend on $I$, and the
hypothesis makes them equal for $I=J$ and $I=\Col$; for instance
$\bigvee_\mu\ell_J(q_\mu)=\bigvee_\mu\ell_{\Col}(q_\mu)$. Hence
$\sigma_J(u,v)=\sigma_{\Col}(u,v)$ for all $u,v\in D_J$, and adding back the
unchanged witnesses gives the identity.
\end{proof}
The lemma applies to any type-constant operands, not only to inputs, and
Lemma~\ref{lem:profile} makes it available at every gate of a computation.
It says that an aggregation cannot distinguish $J$ from $\Col$ unless its
three responses do; how many colors are needed to fix those responses is a
property of the kernel alone.

\begin{definition}[Support and support number]\label{def:width}
A set $K\subseteq\Col$ \emph{supports} the family
$(p_\lambda,q_\lambda)_{\lambda\in\Col}$ if
$(\ell_K,r_K,d_K)=(\ell_{\Col},r_{\Col},d_{\Col})$. The \emph{support number}
$\wid\odot$ of a kernel $\odot$, with respect to the given join, is
the least $c$ such that every finite family, over every finite set of
colors, has a support of size at most $c$; it is $\infty$ if no such $c$
exists. It therefore depends on the join as well as on $\odot$; the join is
always the one fixed with $S$, and Section~\ref{sec:matrix} varies both.
\end{definition}
If $K$ supports a family then so does every $J$ with
$K\subseteq J\subseteq\Col$, since $\ell_K\le\ell_J\le\ell_{\Col}$ pointwise
and likewise for $r$ and $d$. A support may be empty: as $\ell_\varnothing$
and $r_\varnothing$ are identically $0$ and $d_\varnothing=0$, the empty set
supports exactly the families whose three responses to $\Col$ vanish, which
for $\odot=\land$ happens as soon as every $p_\lambda$ and every $q_\lambda$
is $0$. For a general kernel that last inference is not available, since no
algebraic law is assumed for $\odot$ and $0\odot s$ need not be $0$. Either
way an empty support causes no trouble, since it is the retained set $J$ of
colors, never the support, that is required to be nonempty.

For the kernel that Theorems~\ref{main:omega} and~\ref{main:phi} need, two
colors always suffice.
\begin{lemma}[Conjunction has support number two]\label{lem:and-width}
For $S=\bool$, with disjunction as the join and $\odot=\land$, one has
$\wid\land=2$.
\end{lemma}
\begin{proof}
Here $\ell_J(s)=s\land\bigvee_{\lambda\in J}p_\lambda$,
$r_J(s)=s\land\bigvee_{\lambda\in J}q_\lambda$, and
$d_J=\bigvee_{\lambda\in J}(p_\lambda\land q_\lambda)$. So $K$ supports the
family iff it contains a color with $p_\lambda=1$ whenever there is one, a
color with $q_\lambda=1$ whenever there is one, and a color with
$p_\lambda=q_\lambda=1$ whenever there is one. If some color has
$p_\lambda=q_\lambda=1$, it alone supports the family; otherwise one color
with $p_\lambda=1$ and one with $q_\lambda=1$, whenever such colors exist,
together support it. The family
$(p_\lambda,q_\lambda)=(1,0),(0,1)$ on two colors needs both of them, so
$\wid\land\ge2$.
\end{proof}
The \emph{support number} of a computation $\Comp$ is the sum of the
support numbers of its aggregation gates,
\[
 \wid\Comp=\sum_{g\text{ an aggregation gate of }\Comp}\wid{\odot_g},
\]
so a \CoR\ term or relational circuit with $m$ compositions has support
number $2m$.

\subsection{Simultaneous preservation}\label{sec:simultaneous}
As outlined in Section~\ref{ov:wholeterm}, choosing a support at every
gate and taking the union preserves every aggregation of a computation at
once.
\begin{theorem}[Preservation]\label{thm:preservation}
Let $\Comp$ be a computation whose inputs are type-constant arrays on $D$.
There is a set $K\subseteq\Col$ with $|K|\le\wid{\Comp}$ such that for
every nonempty $J$ with $K\subseteq J\subseteq\Col$ and every gate $g$ of
$\Comp$,
\[
 \text{value of $g$ on the inputs restricted to $D_J$}
 \;=\;
 \bigl(\text{value of $g$ on $D$}\bigr)\res J .
\]
In particular the outputs agree. Each aggregation gate is charged once,
however often its value is used.
\end{theorem}
\begin{proof}
If $\wid{\Comp}=\infty$, take $K=\Col$. Otherwise evaluate $\Comp$ on
$D$. By Corollary~\ref{cor:closure}, extended to arbitrary kernels as noted
in Section~\ref{sec:responses}, the value of every gate is type-constant.
For each aggregation gate $g$ with operands $P,Q$ choose a support $K_g$ of
the family $(P(\ctype\lambda),Q(\ctype\lambda))_{\lambda\in\Col}$ with
$|K_g|\le\wid{\odot_g}$, and let $K=\bigcup_gK_g$. Now fix a nonempty
$J\supseteq K$ and proceed by induction along the evaluation order. For the
inputs the claim holds by definition. A free gate commutes with
restriction, so if its operands agree after restriction, so does its value.
At an aggregation gate the operands agree by induction, $J\supseteq K_g$
supports its reference operands, and Lemma~\ref{lem:local} gives the claim.
The supports were chosen on the reference computation, independently of
$J$; this is why a single set $K$ serves every $J\supseteq K$.
\end{proof}
Two remarks are in order. First, the only closure property used is that of the values on
$D$; nothing is assumed about the restricted values. Second, for a \CoR\
term or circuit with $m$ compositions Lemma~\ref{lem:and-width} gives
$|K|\le2m$; complement gates need no separate treatment, since the theorem
asserts an equality (Section~\ref{ov:wholeterm}).

The theorem reduces lower bounds to counting. Let $\varphi$ be a binary
query and let $\M$ interpret its symbols by type-constant relations. A color
$\lambda$ is \emph{detectable} for $\varphi$ if
\[
 \varphi^{\M\res(\Col\setminus\{\lambda\})}(u,v)\ne\varphi^{\M}(u,v)
 \qquad\text{for some }u,v\in D_{\Col\setminus\{\lambda\}},
\]
that is, if deleting the cell of color $\lambda$ changes $\varphi$ at a pair
of surviving vertices.

\begin{corollary}[Detectable colors]\label{cor:detectable}
Let $\Det\subseteq\Col$ be a set of detectable colors for $\varphi$ on $\M$.
Every computation $\Comp$ whose output equals $\varphi$ on $\M$ and on the
$|\Det|$ structures $\M\res(\Col\setminus\{\lambda\})$, $\lambda\in\Det$,
has $\wid{\Comp}\ge|\Det|$. In particular, every \CoR\ term or relational
circuit that agrees with $\varphi$ on these $|\Det|+1$ finite structures has
at least $\lceil|\Det|/2\rceil$ compositions, since that count is an integer.
\end{corollary}
\begin{proof}
If $\wid{\Comp}<|\Det|$, the set $K$ of Theorem~\ref{thm:preservation}
misses some $\lambda\in\Det$. For $J=\Col\setminus\{\lambda\}$, which is
nonempty as $N\ge2$, the output of $\Comp$ on $\M\res J$ is the restriction
of its output on $\M$, whereas $\varphi$ changes at a pair of $D_J$.
\end{proof}
For a non-Boolean $S$, ``the output equals $\varphi$'' means that the output
encodes the truth values of $\varphi$ by two fixed distinct elements of $S$;
intermediate values may be arbitrary. In both interpretations of
Section~\ref{sec:logical} every color is detectable, so $\Det=\Col$ and the
bound is $\lceil|\Col|/2\rceil$; Remark~\ref{rem:iff} in
Appendix~\ref{app:optimality} is the one place where a single color has to
be given up.
